\clearpage\phantomsection\addcontentsline{toc}{chapter}{چکیده}
% SOURCE English_Abstract.txt:20-20 [paragraph]
چکیده فارسی (مطابق نسخه اصلی پایان‌نامه):\par

% SOURCE English_Abstract.txt:21-21 [spacing]


% SOURCE English_Abstract.txt:22-22 [paragraph]
عنوان پایان‌نامه: چارچوب امنیتی هیبریدی برای انتقال تصاویر پزشکی در شبکه اینترنت اشیای پزشکی \lr{(IoMT)} از طریق قطعه‌بندی یادگیری عمیق، سیستم ۵بعدی ابرآشوبی و پنهان‌نگاری آگاه به برجستگی\par

% SOURCE English_Abstract.txt:23-23 [spacing]


% SOURCE English_Abstract.txt:24-24 [paragraph]
چکیده:\par

% SOURCE English_Abstract.txt:25-25 [paragraph]
گسترش روزافزون پزشکی از راه دور، سلامت الکترونیک \lr{(E{-}Health)} و شبکه اینترنت اشیای پزشکی \lr{(IoMT)} امکان تبادل سریع داده‌های بالینی را فراهم ساخته است، اما هم‌زمان انتقال تصاویر پزشکی حساس در بسترهای ارتباطی ناامن، این داده‌ها را در معرض تهدیدات سایبری شدید نظیر دستکاری اطلاعات، شنود غیرمجاز و حملات رمزشناسی قرار داده است. الگوریتم‌های رمزنگاری سنتی عمدتاً از کلیدهای ایستا استفاده می‌کنند که سیستم را در برابر حملات متن‌واضح برگزیده \lr{(CPA)} و معلوم \lr{(KPA)} آسیب‌پذیر می‌سازد، یا اعوجاج‌های دیداری ناخواسته‌ای در نواحی تشخیصی ایجاد می‌نمایند. علاوه بر این، پیچیدگی محاسباتی بالای این الگوریتم‌ها مانع از پیاده‌سازی زمان‌واقعی بر روی تجهیزات کم‌مصرف پردازش لبه می‌گردد.\par

% SOURCE English_Abstract.txt:26-26 [spacing]


% SOURCE English_Abstract.txt:27-27 [paragraph]
جهت حل چالش سه‌گانه متعارض امنیت، کیفیت دیداری و سرعت پردازش، این پایان‌نامه یک چارچوب امنیتی هیبریدی ۴لایه‌ای برای انواع مدالیته‌های تصویربرداری پزشکی (شامل \lr{DRIVE}، \lr{RITE}، \lr{BraTS 2020}، و \lr{COVID{-}19 CXR}) ارائه می‌دهد. در لایه اول، شبکه عصبی عمیق \lr{U{-}Net} ناحیه حساس تشخیصی \lr{(ROI)} را به صورت خودکار قطعه‌بندی کرده و با استخراج گشتاورهای آماری آن و اعمال تابع هَش \lr{SHA{-}256}، کلید ۲۵۶ بیتی پویا ($K_{\mathrm{hash}}$) و کاملاً تصویر{-}پایه تولید می‌کند. در لایه دوم، کلید $K_{\mathrm{hash}}$ متغیرهای اولیه سیستم ۵بعدی ابرآشوبی با دو توان مثبت لیپانوف ($L_1 = 1.24,\allowbreak{} L_2 = 0.85$) را مقداردهی نموده و فضای کلید عظیم $2^{512}$ را ایجاد می‌نماید. در لایه سوم، ترکیب جایگشت متقاطع زیگ‌زاگی و انتشار پویای \lr{DNA}، همبستگی پیکسل‌های مجاور را به کمتر از $0.0008$ تنزل داده و آنتروپی شانون را به حد آرمانی $7.9985$ بیت می‌رساند.\par

% SOURCE English_Abstract.txt:28-28 [spacing]


% SOURCE English_Abstract.txt:29-29 [paragraph]
در لایه چهارم، فراداده هویتی بیمار پس از فشرده‌سازی \lr{zlib}، صرفاً در ۲ بیت کم‌ارزش پس‌زمینه کم‌برجستگی جای‌دهی می‌شود. این استراتژی صیانت ۱۰۰ درصدی از ناحیه تشخیصی را تضمین نموده ($\mathrm{PSNR}_{\mathrm{ROI}} = \infty$) و کیفیت استگو کل تصویر را به $\mathrm{PSNR} = 52.41\mathrm{~dB}$ و $\mathrm{SSIM} = 0.9982$ می‌رساند. ارزیابی‌های پایداری اثبات می‌کنند که سیستم تحت ۵۰\SourceGlyph{066A} قطع داده، بافت تشخیصی را با کیفیت مناسب ($\mathrm{PSNR} > 32.0\mathrm{~dB}$) بازیابی کرده و تمام ۱۵ آزمون استاندارد تصادفی‌سازی \lr{NIST SP800{-}22} را با موفقیت پاس می‌نماید ($p \ge 0.0125$). ارزیابی‌های سخت‌افزاری بر روی بستر \lr{NVIDIA Jetson TX2} زمان اجرای کل ۲.۸۵ ثانیه را به ثبت رسانده که قابلیت پیاده‌سازی زمان‌واقعی سیستم در شبکه \lr{IoMT} را اثبات می‌کند.\par

% SOURCE English_Abstract.txt:30-30 [spacing]


واژه‌های کلیدی: امنیت تصاویر پزشکی، یادگیری عمیق \lr{U{-}Net}، آشوب پنج‌بعدی، رمزنگاری \lr{DNA}، پنهان‌نگاری آگاه به برجستگی، اینترنت اشیای پزشکی و پردازش لبه.\par
